\section{Building and End-to-end Platform}

\noindent\textbf{Integrating a Reconfigurable LLC Partitioning Scheme}
\label{sec:dynamiccache}
We build a fine-grain set-partitioning scheme similar to \cite{townley2022composable,saileshwar2021bespoke} to support our 64 memory partitions.
Thanks to some of Citadel design choices, we are able to circumvent challenges in end-to-end integration ignored in previous work.
For instance, our memory region IDs are directly accessible through physical memory address (Section \ref{sec:hardwareisolation}) which makes it possible to support shared memory between cores without extending every structure or buses in the memory hierarchy to support security domain ID.
Additionally, we can leverage hardware software co-desing to ensure the LLC will not be put in an incoherent state during reconfiguration.
A request for reconfiguration always goes through the SM which will check 
that the mapping is correct and fits within the LLC.
The SM will check that no enclaves are running and preempt other cores to ensure no requests will be sent to the affected memory regions during the operation.
Implementation details can be found in Appendix \ref{app:dynamiccache}.
To our knowledge, we are the first to implement such a scheme in RTL and to integrate it in a system.
\\

\noindent\textbf{Fine-grained LLC Flushing}
We also implement a hardware-software fine-grain flushing mechanism.
We add a simple piece of logic connected to the cache hierarchy at the same level than the DRAM, which always sends back zeroes.
It is mapped to a physical address space the size of the DRAM making it possible to easily build and access eviction sets in software.
This means that flushing a small LLC partition is much faster than flushing a big one.
Implementation details can be found in Appendix \ref{app:llcflush}.
\\

\noindent\textbf{Building the Software Infrastructure}
Our software architecture, illustrated early in the paper in Figure \ref{fig:computing_stack}, comprises several key components. 
We implement the SM at RISC-V's highest privilege level using assembler and C, alongside a \emph{Mini-SM} - a downsized, enclave-private copy of the global SM.
A Linux Kernel module enables user programs to request SM calls and set up enclaves, with Linux managing resources.
Additionally, we develop a Secure Bootloader and an end-to-end attestation mechanism, enabling unique enclave identification, shared secret establishment, and encrypted private data exchange.
Implementation details regarding our entire software infrastructure can be found in Appendix \ref{app:softwareinfrastructure}.
